\section{Routing interpretation：离散 gate 提供证据，但不自动给出语义}

Dense network 的 activation 位于高维连续空间，很难直接解释；SMoE router 至少给出离散 expert IDs 和 gating weights，让研究者可以统计 token 到 expert 的分配。本章沿课程的三个分析层次展开：domain histogram、consecutive-token persistence、token-level map；最后再讨论社区 expert ablation 和“expert 不是 domain”的结论。

\subsection{为什么 discrete gate 看起来更可解释}

Router 每层为 token 产生一个小集合 $S_k(x)$，这比数千维 activation 更容易计数。研究者可以问：不同数据 domain 是否偏向不同 experts？相邻 token 是否保持同一路由？删除某个 expert 会怎样？但这些问题只产生 observable correlations，解释仍需要对照实验。

\lecturefigure{slide-19-discrete-gating.jpg}{Discrete gating signals 为 SMoE interpretation 提供新观测量}{Stanford Online 官方录像 00:20:19}

对 domain $d$、layer $\ell$、expert $i$，routing frequency 可写成
\[
p_{\ell}(i\mid d)=\frac{\#\{x\in d:i\in S_k^{(\ell)}(x)\}}{k\,|d|}.
\]
若完全均匀，$p_{\ell}(i\mid d)\approx 1/E$。偏离均匀只是“该 domain 的 token 更常选这个 expert”，不是“这个 expert 的语义就是该 domain”。

\begin{warningbox}{Attention weight 与 router ID 都不是因果解释}
可视化告诉我们模型看了什么或选了谁，但不等于该组件对输出具有唯一因果作用。需要 intervention、ablation、counterfactual routing 或 retraining 才能更接近因果结论。
\end{warningbox}

\teachervoice{00:19:15--00:20:24，讲者认为 discrete gate 比 dense activation 提供了“incredible opportunity”，但他的提问仍是“can we make some sense out of this”，而非宣布可解释性已经解决。}

\subsection{Domain histogram：中层 spike 有趣，但证据仍弱}

实验把 The Pile 的不同 validation subsets 输入模型，统计 layer 0、15、31 的 expert selection。浅层与深层大体接近均匀，中层某个 expert 对 DeepMind Mathematics/GitHub 有较高选择率。读图时先看 $1/E=12.5\%$ baseline，再比较层间变化，而不是给颜色直接命名。

\lecturefigure{slide-20-domain-specialisation.jpg}{The Pile 不同 domains 在浅层、中层和深层的 expert routing 分布}{Stanford Online 官方录像 00:22:39}

\begin{knowledgebox}{读图：三层分别可能看见什么}
Layer 0 接近 raw token，路由可能受字符、标点和局部句法影响；中层包含更多组合表示，因此出现 domain-correlated spike；最后一层接近 decoding objective，分布又可能趋于共享输出需求。这个叙述是可检验 hypothesis，不是由一张图唯一确定。
\end{knowledgebox}

\teachervoice{00:20:24--00:22:28，Albert 说 layer 15 expert 3 的 math/code spike “interesting”，随后马上补充 highly speculative、cannot conclude much。这个自我限制必须保留。}

\subsection{Consecutive tokens：routing persistence 在中层更强}

Domain histogram 只告诉我们不同数据子集的总体选择比例，无法判断 router 是否在一个局部片段内保持决策。第二个实验因此把问题从“什么 domain 更常选谁”改成“相邻 token 是否连续选择同一 expert”。读图时先确认 top-1 与 top-2 两种事件的随机 baseline，再比较浅层、中层和深层，而不要把连续性直接解释为主题专门化。若 top-1 在八个 experts 间独立均匀，随机 baseline 为 $1/8=12.5\%$；若比较“下一 token 是否把该 expert 放在 top-2”，baseline 会更高。图中 layer 15 的 persistence 明显高于随机，layer 0 较弱，layer 31 又回落。

\lecturefigure{slide-21-consecutive-tokens.jpg}{相邻 token 继续选择同一 expert 的频率}{Stanford Online 官方录像 00:25:13}

定义 top-1 persistence
\[
P_{\ell}^{(1)}=\Pr\!\left(e_{t+1}^{(1)}=e_t^{(1)}\right),
\]
以及 set persistence
\[
P_{\ell}^{(2)}=\Pr\!\left(e_t^{(1)}\in S_2(x_{t+1})\right).
\]
它们衡量局部路由连续性，却不能说明连续 token 共享“主题专家”；subword、syntax、position 或 hidden-state smoothness 都可能造成 persistence。

\teachervoice{00:22:46--00:25:18，讲者逐层解释 random baseline 与 observed frequency：layer 15 约为最强，最后一层仍高于随机但有所回落，并明确说需要更细分析。}

\subsection{Token map：颜色看起来有模式，语义却不整齐}

第三个实验把代码、算术和选择题 token 按 expert 着色。某些数字或局部片段出现相似颜色，但整体没有“代码全部进 expert A、数学全部进 expert B”的清晰划分。图像最重要的教学作用，是让读者亲眼看到 expert routing 的碎片化，而不是制造一个漂亮 taxonomy。

\lecturefigure{slide-22-token-assignment.jpg}{代码、算术与问答样本的 token-level expert assignment}{Stanford Online 官方录像 00:26:13}

\begin{knowledgebox}{读图：先看局部连续，再看跨样本一致性}
第一步找同一序列中的连续颜色块；第二步比较同类 token 在不同样本是否保持颜色；第三步跨 layer 比较。只有多样本、跨层、跨随机种子都稳定，才值得尝试给 expert 语义命名。
\end{knowledgebox}

\teachervoice{00:25:18--00:26:19，Albert 观察到一些 digits 被送到同一 expert，但总体“doesn't seem to be much specialization”。课程没有从 token map 推出明确 domain taxonomy。}

\subsection{Myth 4：为什么纯 domain experts 还可能降低利用率}

假设只有两个 experts 专门处理 code，那么长代码序列会把大量 token 压到这两个 experts，其余六个空闲；这既产生 load imbalance，也浪费并行 capacity。更合理的 specialization 可能是跨 domain 共享的 latent features，使所有 experts 在不同输入上都能被利用。

\lecturefigure{slide-23-myth-domain-specialisation.jpg}{Myth 4：最优 experts 不一定按人类 domains 分工}{Stanford Online 官方录像 00:27:15}

若 domain $d$ 的 token 全部路由到子集 $A_d$，利用率上界近似为
\[
U_d=\frac{|A_d|}{E}.
\]
例如 $|A_d|=2,E=8$ 时，单一 domain batch 最多充分使用四分之一 experts。现实 router 可 top-two 混合并跨 feature 分工，但公式揭示“干净 domain 专家”与高利用率之间的张力。

\begin{importantbox}{Specialization 的更稳妥定义}
不要先指定 expert 名称，再找例子；应测量路由分布、representation probe、causal ablation 与 load behavior，再描述某个 expert 对哪些 features 敏感。Feature 可以跨代码、数学和自然语言同时出现。
\end{importantbox}

\teachervoice{00:26:20--00:27:22，讲者用 coding experts 的反例说明：若代码 token 只进两个 experts，其他 experts 就站着不工作；语言太复杂，不应把 specialization 简化成高层 domains。}

\subsection{Treasure hunt：开放权重让社区快速做 intervention}

模型发布约一天后，社区尝试移除特定 expert，发现删除某一编号会让 MMLU 大幅崩溃，并做成“一个人在干活、其他人在围观”的 meme。这个结果很醒目，但它可能反映 router calibration、layer-specific bottleneck、implementation 或 distributed ablation 方法；不能直接证明一个 expert 存了全部知识。

\lecturefigure{slide-24-treasure-hunt.jpg}{社区 expert ablation、MMLU 崩溃与开源 meme}{Stanford Online 官方录像 00:28:31}

定义删除 expert $i$ 后的 metric change
\[
\Delta_i=m(\theta_{-i})-m(\theta),
\]
其中 $\theta_{-i}$ 表示禁用某 expert 的干预模型。大的负 $\Delta_i$ 证明该干预影响输出，却不能区分 expert 本身知识、router re-normalization 或 downstream distribution shift。

\begin{warningbox}{Ablation 设计必须记录四件事}
删除的是一个 layer 的 expert 还是所有层同编号？Gates 是否重新归一化？Overflow token 如何处理？Metric 是否对输出格式敏感？不记录这些，社区结果只能是线索。
\end{warningbox}

\teachervoice{00:27:22--00:28:38，Albert 把这个 24 小时内出现的社区发现当作“why we love open source”的例子；价值在于快速 intervention，不在于 meme 已经解释模型。}

\subsection{最终解释议程：寻找 latent feature subspaces}

课程最后没有给 experts 贴标签，而是提出更难的问题：experts 可能编码人类概念的线性组合，或对某个 feature subspace 敏感。未来研究需要把 routing stats 与 representation analysis、causal patching、expert swapping 和 controlled retraining 结合起来。

\lecturefigure{slide-25-interpret-moe-decisions.jpg}{开放问题：experts 真正学习了哪些 features}{Stanford Online 官方录像 00:29:11}

\begin{importantbox}{从可视化到解释的证据阶梯}
Routing histogram $\rightarrow$ token map $\rightarrow$ controlled ablation $\rightarrow$ counterfactual rerouting $\rightarrow$ retraining/feature recovery。越往后越接近因果，但成本也越高。单张 histogram 只能位于阶梯起点。
\end{importantbox}

\teachervoice{00:28:38--00:29:30，讲者认为 experts 可能捕获与人类概念不同的 features，甚至是多个概念的线性组合；真正问题是怎样恢复其 subspace。}

\subsection{本章小结}

Discrete routing 确实给 interpretability 增加了可观测变量：domain frequency、routing persistence、token assignment 和 ablation sensitivity。但所有结果都指向同一谨慎结论：experts 具有结构化行为，却未按简单人类 domain 分工。最有价值的研究方向是把离散 gate 与因果干预、representation geometry 和系统负载联合起来。

